% FUTURE WORK — drafted 2026-09-08. Prose numbers for now; macros later.

\section{Future work}
\label{sec:future_work}

\subsection{The other three languages}

The corpus is Kinyarwanda only, and the three other languages are at three
different stages. None of them contributes a row to any result in this paper.

\paragraph{English and French are drafted but unverified.} Both exist as
256-row review briefs on the same 128-concept spine, carrying candidate
phrasings with per-row provenance and review verdicts. Every candidate is
machine-drafted. No English speaker and no French speaker has ruled on them,
and the French brief records explicitly that no French speaker has seen any of
it. They are not a corpus and are not treated as one.

The two arms did produce a result worth keeping: built independently against
the same spine, they lifted the same eleven Kinyarwanda holds on the same
reasoning --- each being a block on a Kinyarwanda \emph{word} rather than on
the concept --- and they added the same hold on a concept-level clinical
question. Two arms agreeing independently is the strongest evidence available
that those rulings are about concepts rather than about one language.

\paragraph{Swahili is out for authoring.} A Kiswahili speaker has agreed to
author the arm, and the brief they received is an authoring brief rather than a
review brief: it contains no Swahili at all. This is deliberate and follows the
Kinyarwanda finding --- a fluent draft anchors a speaker to its errors, which is
why the Kinyarwanda arm was re-authored from scratch rather than corrected. The
brief carries the English gloss, the person split, the applicability and
relation rulings and the open clinical questions, and an empty column for the
speaker's own phrasing. A machine-translated Swahili corpus exists in the frozen
v1 vocabulary and is deliberately withheld from the speaker; it can be compared
against the authored phrases afterwards, which is the only order in which the
comparison means anything.

Two gaps in the Swahili arm are recorded and not yet closed. No Swahili relation
terms exist anywhere in the project, so every third-person placeholder currently
has nothing to substitute; and the frame slots --- greetings, time expressions,
closing lines --- have not been authored either.

\subsection{The evaluation base}

The limitation in Section~\ref{sec:limitations} is the one that most constrains
what can be claimed, and it is a property of the corpus rather than of the
model. Enlarging the number of distinct phrases per class, rather than the
number of rows, is the only thing that would make a recall target of 0.95
measurable at better than quarter granularity.

\subsection{The unresolved boundary}

Every configuration trained separated ROUTINE without error and none resolved
the CRITICAL/URGENT boundary. The next diagnostic is per-sentence rather than
per-configuration: which of the held-out CRITICAL sentences is read as URGENT,
and whether it is the same sentence across models. That distinguishes an
ambiguous concept from a mislabelled one from an artefact of the split, and it
is one inference pass rather than another training run.

\subsection{Thresholds that are still not sourced}

Two numbers in the acceptance gate need an authority this project does not
have. The CRITICAL recall threshold of 0.95 is inherited and its source is not
verifiable in this repository. The margin by which CRITICAL precision must
exceed the degenerate ceiling is a judgement about tolerable over-triage in a
clinic, and no clinician has made it. Both are reported as placeholders.
